\section{Rozmaitości modelowane na przestrzeniach Banacha}
\label{sec:fsg-banach-manifolds}

\subsection{Modele lokalne i zgodność map}

Przestrzeń Banacha jest naturalnym modelem lokalnym przestrzeni
funkcji. W odróżnieniu od przestrzeni euklidesowej może mieć
nieskończony wymiar, a jej domknięte podprzestrzenie nie muszą
mieć domkniętych dopełnień. Dlatego w dalszej teorii istotne
będzie nie tylko istnienie pochodnej, lecz również struktura
jej jądra i obrazu.

Ustalamy rzeczywistą przestrzeń Banacha $E$ oraz
$k\in\{1,2,\ldots,\infty\}$. Klasę $C^k$ odwzorowań między
otwartymi podzbiorami przestrzeni Banacha rozumiemy w sensie
Definicji~\ref{def:ck15}.

\begin{definicja}[Mapa i atlas]
\label{def:fsg-chart}
Niech $M$ będzie przestrzenią topologiczną Hausdorffa.
Mapą o modelu $E$ nazywamy parę $(U,\varphi)$, gdzie
$U\subseteq M$ jest zbiorem otwartym, a
$\varphi:U\to V$ jest homeomorfizmem na otwarty podzbiór
$V\subseteq E$. Dwie mapy $(U,\varphi)$ i $(W,\psi)$
są zgodne klasy $C^k$, jeżeli odwzorowania przejścia
\[
 \psi\circ\varphi^{-1}:\varphi(U\cap W)\longrightarrow
 \psi(U\cap W)
\]
oraz jego odwrotność są klasy $C^k$.
Atlas klasy $C^k$ jest rodziną parami zgodnych map,
których dziedziny pokrywają $M$.
Rozmaitością Banacha klasy $C^k$ nazywamy przestrzeń
$M$ wraz z maksymalnym takim atlasem.
Nie zakładamy w tej definicji drugiego aksjomatu przeliczalności
ani parazwartości; jeżeli będą potrzebne, wskażemy je osobno.
\end{definicja}

\begin{stwierdzenie}[Maksymalne rozszerzenie atlasu]
\label{prop:fsg-maximal-atlas}
Każdy atlas klasy $C^k$ wyznacza dokładnie jeden
maksymalny atlas klasy $C^k$, który go zawiera.
\end{stwierdzenie}
\begin{proof}
Do wyjściowego atlasu $\mathcal A$ dołączamy wszystkie mapy
zgodne z każdą mapą z $\mathcal A$. Aby sprawdzić zgodność
dwóch dołączonych map $(U,\varphi)$, $(W,\psi)$,
ustalmy $p\in U\cap W$. Wybieramy mapę $(O,\theta)\in\mathcal A$
taką, że $p\in O$. Na otwartym otoczeniu $\varphi(U\cap W\cap O)$
punktu $\varphi(p)$ mamy
\[
 \psi\circ\varphi^{-1}
 =(\psi\circ\theta^{-1})\circ(\theta\circ\varphi^{-1}).
\]
Oba czynniki są klasy $C^k$, więc ich złożenie też jest tej
klasy. Analogicznie postępujemy z przejściem odwrotnym.
Ponieważ własność $C^k$ jest lokalna, otrzymujemy zgodność
na całym przecięciu. Maksymalność i jednoznaczność wynikają
z tego, że każda mapa atlasu zawierającego $\mathcal A$
musi być zgodna z każdą mapą z $\mathcal A$.
\end{proof}

\begin{przyklad}[Otwarte podzbiory przestrzeni funkcyjnych]
\label{ex:fsg-positive-functions}
Niech $K$ będzie niepustą zwartą przestrzenią Hausdorffa,
a $E=C(K;\R)$ z normą supremum. Zbiór
\[
 \mathcal P=\{f\in E:f(x)>0\text{ dla każdego }x\in K\}
\]
jest otwarty. Istotnie, dla $f\in\mathcal P$ mamy
$m=\min_{x\in K}f(x)>0$. Jeśli $\|g-f\|_\infty<m/2$,
to $g(x)>m/2$ dla każdego $x\in K$. Mapa tożsamościowa
na $\mathcal P$ wyznacza atlas klasy $C^\infty$.
Jest to rozmaitość Banacha, mimo że jej punkty są funkcjami,
a nie wektorami o skończonej liczbie współrzędnych.
\end{przyklad}

\subsection{Odwzorowania i przestrzenie styczne}

\begin{definicja}[Odwzorowanie klasy $C^r$]
\label{def:fsg-smooth-map}
Niech $M,N$ będą rozmaitościami Banacha klasy $C^k$,
modelowanymi odpowiednio na $E,F$, i niech $1\le r\le k$.
Ciągłe odwzorowanie $f:M\to N$ jest klasy $C^r$, jeżeli
dla każdej pary map $(U,\varphi)$, $(W,\psi)$
odwzorowanie
$\psi\circ f\circ\varphi^{-1}$ jest klasy $C^r$
na otwartym zbiorze $\varphi(U\cap f^{-1}(W))$.
\end{definicja}

Wystarczy sprawdzić ten warunek lokalnie przy jednej parze
map wokół każdego punktu i jego obrazu: dla innych map
lokalny zapis zmienia się przez złożenie z odwzorowaniami
przejścia klasy $C^k$. Reguła łańcuchowa
(Twierdzenie~\ref{thm:chain14}) i jej iteracje zapewniają
zachowanie klasy $C^r$.

\begin{definicja}[Wektor styczny jako klasa krzywych]
\label{def:fsg-tangent}
Niech $p\in M$. Rozważamy krzywe klasy $C^1$
$\gamma:(-\varepsilon,\varepsilon)\to M$ z $\gamma(0)=p$.
Dwie takie krzywe uznajemy za równoważne w punkcie $p$,
jeśli dla pewnej mapy $(U,\varphi)$ wokół $p$
\[
 (\varphi\circ\gamma_1)'(0)
 =(\varphi\circ\gamma_2)'(0).
\]
Przestrzeń styczna $T_pM$ jest zbiorem klas równoważności
tych krzywych.
\end{definicja}

\begin{twierdzenie}[Niezależność od mapy i struktura liniowa]
\label{thm:fsg-tangent-chart}
Relacja z Definicji~\ref{def:fsg-tangent} nie zależy od
wyboru mapy. Dla każdej mapy $(U,\varphi)$ wokół $p$
odwzorowanie
\[
 I_{\varphi,p}:T_pM\to E,\qquad
 I_{\varphi,p}([\gamma])=(\varphi\circ\gamma)'(0)
\]
jest bijekcją. Przeniesienie struktury przestrzeni Banacha
przez tę bijekcję daje strukturę liniową i topologię
niezależne od mapy, choć konkretna norma może zależeć od mapy.
\end{twierdzenie}
\begin{proof}
Dla drugiej mapy $(W,\psi)$ reguła łańcuchowa daje
\[
 (\psi\circ\gamma)'(0)
 =D(\psi\circ\varphi^{-1})(\varphi(p))
   (\varphi\circ\gamma)'(0).
\]
Pochodna przejścia jest ograniczonym izomorfizmem liniowym:
reguła łańcuchowa dla przejścia i jego odwrotności daje
obie tożsamości operatorowe. Równość prędkości w jednej
mapie jest więc równoważna równości w drugiej, co również
dowodzi, że zdefiniowano relację równoważności.
Wstrzykiwanie $I_{\varphi,p}$ wynika z definicji klas.
Dla $v\in E$ wybieramy $\varepsilon>0$ tak, by
$\varphi(p)+tv\in\varphi(U)$ dla $|t|<\varepsilon$.
Krzywa $\gamma(t)=\varphi^{-1}(\varphi(p)+tv)$ ma
współrzędną prędkość $v$, co daje surjekcję.
Zmiana mapy odpowiada ograniczonemu izomorfizmowi liniowemu,
zatem przeniesione działania liniowe są identyczne, a normy
równoważne. Zupełność wynika z zupełności $E$.
\end{proof}

\begin{twierdzenie}[Różniczka odwzorowania rozmaitości]
\label{thm:fsg-tangent-map}
Dla odwzorowania $f:M\to N$ klasy $C^1$ różniczka
\[
 T_pf:T_pM\to T_{f(p)}N,\qquad [\gamma]\mapsto[f\circ\gamma]
\]
jest dobrze określonym ograniczonym operatorem liniowym.
W mapach $\varphi,\psi$ wokół $p,f(p)$ zachodzi
\[
 I_{\psi,f(p)}\circ T_pf\circ I_{\varphi,p}^{-1}
 =D(\psi\circ f\circ\varphi^{-1})(\varphi(p)).
\]
Jeżeli $g:N\to P$ jest klasy $C^1$, to
$T_p(g\circ f)=T_{f(p)}g\circ T_pf$.
\end{twierdzenie}
\begin{proof}
Współrzędna prędkość $f\circ\gamma$ jest, z reguły
łańcuchowej, obrazem prędkości $\varphi\circ\gamma$
przez operator po prawej stronie równości. Zależy więc
wyłącznie od klasy $[\gamma]$. Wzór współrzędny
daje liniowość i ograniczoność. Ostatnia równość wynika
z $(g\circ f)\circ\gamma=g\circ(f\circ\gamma)$.
\end{proof}

\begin{przyklad}[Różniczka funkcjonału całkowego]
\label{ex:fsg-integral-tangent}
Na rozmaitości $\mathcal P\subset C(K;\R)$ z
Przykładu~\ref{ex:fsg-positive-functions} niech $\mu$
będzie skończoną dodatnią miarą borelowską na $K$.
Funkcjonał $J(f)=\int_K f(x)^2\,\mathrm d\mu(x)$
ma różniczkę
\[
 T_fJ(h)=2\int_K f(x)h(x)\,\mathrm d\mu(x),
 \qquad h\in T_f\mathcal P\simeq C(K;\R).
\]
Istotnie, reszta po odjęciu części liniowej jest równa
$\int_K h^2\,\mathrm d\mu$, a jej moduł nie przekracza
$\mu(K)\|h\|_\infty^2$. Operator liniowy ma normę
nie większą niż $2\mu(K)\|f\|_\infty$.
Ponadto zmiana tego operatora przy zmianie $f$ na $g$
ma normę nie większą niż $2\mu(K)\|f-g\|_\infty$,
więc $J$ jest klasy $C^1$.
\end{przyklad}
